% BHU PYQ -- Manually transcribed & error-corrected
% Paper: BCH-213 : Specialised Accounts - I
% Exam: B.Com. (Hons.) Semester III Examination, 2023-24

\documentclass[11pt,a4paper]{article}
\usepackage[a4paper,top=2.5cm,bottom=2.5cm,left=2.8cm,right=2.8cm,headheight=14pt]{geometry}
\usepackage{amsmath,amssymb}
\usepackage[T1]{fontenc}
\usepackage[utf8]{inputenc}
\usepackage{lmodern,microtype,fancyhdr,enumitem,hyperref,lastpage,parskip}

\hypersetup{
    colorlinks=true,
    urlcolor=black,
    linkcolor=black,
    pdftitle={BCH-213: Specialised Accounts - I -- BHU B.Com. (Hons.) Sem III 2023-24}
}

\pagestyle{fancy}
\fancyhf{}
\renewcommand{\headrulewidth}{0.4pt}
\renewcommand{\footrulewidth}{0.4pt}
\fancyhead[L]{\small   \textit{Banaras Hindu University}}
\fancyhead[R]{\small   \textit{BCH-213: Specialised Accounts - I}}
\fancyfoot[C]{\small Page  \thepage\ of \pageref{LastPage}}


\newlist{parts}{enumerate}{2}
\setlist[parts,1]{label=(\roman*),leftmargin=*,itemsep=5pt,topsep=3pt}
\setlist[parts,2]{label=(\alph*),leftmargin=*,itemsep=3pt,topsep=2pt}

\newcommand{\pts}[1]{\hfill   \textit{[#1]}}


\begin{document}


\begin{center}
  {\large\bfseries BANARAS HINDU UNIVERSITY}\\[2pt]
  {\normalsize B.Com. (Hons.) --- Semester III Examination, 2023-24}\\[6pt]
  \rule{0.8\linewidth}{0.5pt}\\[6pt]
  {\large\bfseries Paper No. BCH-213: Specialised Accounts - I}\\[4pt]
  {\normalsize Time: 3 Hours\hfill Maximum Marks: 70}
\end{center}

\rule{\linewidth}{0.4pt}

\smallskip



\noindent   \textit{Instructions: Answer all questions. All questions carry equal marks (14 marks each).}

\bigskip




\noindent   \textbf{Question 1.}
Show what date-wise entries would be passed by the Allahabad Branch, Kanpur Branch and the Head Office to record the following transactions in their books, assuming that books are closed on 31st March:

\begin{parts}
  \item Goods amounting Rs. 1,000 transferred from Kanpur Branch to Allahabad Branch on 15th March, 2023 under the instructions from Head Office.
  \item Under the instructions of H. O. Allahabad Branch transferred its Debtor's A/c of Rs. 800 to Kanpur Branch on 16th March, 2023, as the debtor has shifted his business from Allahabad to Kanpur.
  \item Depreciation amounting to Rs. 4,000 on fixed assets of Allahabad Branch and Rs. 500 on fixed assets of Kanpur Branch is to be charged when such accounts are opened in the Head Office books.
  \item A remittance of Rs. 3,000 made by the Kanpur Branch to Head Office on 25th March, 2023 and received by the Head Office on 5th April, 2023.
  \item Goods amounting Rs. 5,000 sent by the Head Office to the Allahabad Branch on 20th March, 2023 and received by the latter on 2nd April, 2023.
  \item The Allahabad Branch collected Rs. 2,000 from Kolkata customer of the Head Office on 23rd March, 2023.
  \item The Allahabad Branch paid Rs. 12,000 for machinery purchased by the Head Office in Allahabad on 26th March, 2023.
  \item On 1st March, 2023 the H. O. had sent goods of Rs. 600 to Allahabad Branch which was stolen in transit, and the Insurance company admitted a claim of Rs. 200 only.
  \item On 18th March, 2023, Kolkata Branch returned goods of Rs. 300 to H. O. which was destroyed by fire in transit. The goods were uninsured.
\end{parts}\pts{14}

\centerline{   \textbf{OR}}

From the information given below, prepare the following in the books of H. O.:

\begin{parts}
  \item Journal entries to incorporate the Trial Balance of the Branch:
  
\begin{parts}
    \item Under Detailed Method
    \item Under Concise Method
  \end{parts}
  \item Branch Trading and P \& L A/c and
  \item Branch A/c
\end{parts}


\begin{center}
   \textbf{Branch Trial Balance}\\
   \textbf{(As on 31st March, 2023)}
\smallskip

\small

\begin{tabular}{|l|r|r|}
\hline
   \textbf{Particulars} &   \textbf{Debit (Rs.)} &   \textbf{Credit (Rs.)} \\
\hline
Opening Stock & 7,270 & -- \\
Purchases & 17,350 & -- \\
Commission Received & -- & 200 \\
Goods from H. O. & 8,100 & -- \\
Creditors & -- & 2,340 \\
Selling Expenses & 4,260 & -- \\
Administrative Expenses & 2,040 & -- \\
Sundry Expenses & 1,380 & -- \\
Sales & -- & 38,000 \\
H. O. A/c & -- & 7,560 \\
Petty Cash & 50 & -- \\
Cash at Bank & 1,250 & -- \\
Debtors & 6,400 & -- \\
\hline
   \textbf{Total} &   \textbf{48,100} &   \textbf{48,100} \\
\hline
\end{tabular}
\end{center}




\noindent   \textbf{Adjustments / Additional Information:}

\begin{parts}
  \item Closing stock on 31st March, 2023 is Rs. 8,200.
  \item The Branch A/c in the Head Office books showed a debit balance of Rs. 8,440, while the Goods sent to Branch A/c showed a credit balance of Rs. 8,980, by the closing date.
  \item A provision for doubtful debts is to be created at 1.5\% on Debtors.
\end{parts}\pts{14}

\medskip\rule{\linewidth}{0.2pt}

\pagebreak




\noindent   \textbf{Question 2.}
Cash and hire-purchase prices of a machine are Rs. 30,000 and Rs. 40,000 respectively. This machine was sold on hire-purchase system on 1st April, 2020 and payment is received in four six-monthly instalments of Rs. 10,000 each. It was agreed that the machine will be maintained free of charge for a period of two years. The actual expenses of maintenance were Rs. 1,000 and Rs. 2,500 for the first and second year, respectively, though past experience showed that they are Rs. 5,000 of which Rs. 700 are for the first year.

Pass the necessary Journal Entries in the books of the hire-vendor and open the Maintenance Suspense A/c. Accounts are closed on 31st March.\pts{14}

\centerline{   \textbf{OR}}

Explain the concept of Hire-Purchase and Instalment Payment System. What journal entries are passed in the books of buyer and seller when goods are sold on hire-purchase system?\pts{14}

\medskip\rule{\linewidth}{0.2pt}




\noindent   \textbf{Question 3.}
Prof. Ram wrote a book on Accountancy and got it published with 'Navyug' on the terms that royalties will be 15\% on the published price of the copies sold with a minimum of Rs. 1,00,000 per year. Prof. Ram gave an undertaking to revise the book when requested by the publisher and to pay Rs. 5,000 to 'Navyug' per month for every month of delay after six months of the request made by 'Navyug'. In the event of delay, the condition of the minimum amount of Rs. 1,00,000 payable to Prof. Ram was not to apply. The arrangement was to last for 10 years. The number of copies sold and the published price were as under:


\begin{center}
\small

\begin{tabular}{|l|c|c|c|c|c|}
\hline
   \textbf{Year} &   \textbf{2018} &   \textbf{2019} &   \textbf{2020} &   \textbf{2021} &   \textbf{2022} \\
\hline
No. of copies sold & 2,000 & 4,000 & 5,000 & 2,000 & 5,000 \\
Published price per book (Rs.) & 200 & 200 & 250 & 250 & 300 \\
\hline
\end{tabular}
\end{center}

At the end of the third year, Prof. Ram was requested to revise the book, and the revised manuscript reached 'Navyug' after 10 months of the request.

Record journal entries in the books of both the parties.\pts{14}

\centerline{   \textbf{OR}}

What journal entries are passed by lessee and landlord in connection with royalty? State the concept and functions of Rat-Hole Mining.\pts{14}

\medskip\rule{\linewidth}{0.2pt}

\pagebreak




\noindent   \textbf{Question 4.}
From the following balances extracted from the books of B. B. K. Lal, prepare Departmental Trading and Profit \& Loss Account for the year ended 31st March, 2022 and Balance Sheet as on that date:


\begin{center}
\small

\begin{tabular}{|l|r|r|}
\hline
   \textbf{Particulars} &   \textbf{Debit (Rs.)} &   \textbf{Credit (Rs.)} \\
\hline
Capital & -- & 1,00,000 \\
Drawings & 16,000 & -- \\
Buildings & 5,00,000 & -- \\
Furniture (Deptt. A) & 50,000 & -- \\
Opening Stock (Deptt. A) & 6,000 & -- \\
Opening Stock (Deptt. B) & 8,000 & -- \\
Purchases (Deptt. A) & 1,00,000 & -- \\
Purchases (Deptt. B) & 1,50,000 & -- \\
Sales (Deptt. A) & -- & 4,00,000 \\
Sales (Deptt. B) & -- & 6,00,000 \\
Carriage Inward & 10,000 & -- \\
Electricity Expenses & 21,000 & -- \\
Salaries & 90,000 & -- \\
Rent & 80,000 & -- \\
Commission on Sales & 60,000 & -- \\
Audit Fees & 10,000 & -- \\
Sundry Debtors & 50,000 & -- \\
Sundry Creditors & -- & 65,000 \\
Cash at Bank & 14,000 & -- \\
\hline
   \textbf{Total} &   \textbf{11,65,000} &   \textbf{11,65,000} \\
\hline
\end{tabular}
\end{center}




\noindent   \textbf{Additional Information:}

\begin{parts}
  \item Depreciate building by 5\% p.a. and furniture by 10\% p.a.
  \item Floor area occupied by each department is equal.
  \item Electricity consumed is in the ratio of 1:2.
  \item Audit fees be apportioned on the basis of turnover.
  \item Number of employees is in the ratio of 2:1.
  \item Closing Stock: Deptt. A = Rs. 26,000 and Deptt. B = Rs. 52,000.
\end{parts}\pts{14}

\centerline{   \textbf{OR}}

Explain objectives of preparing Departmental Accounts. Discuss the various bases of allocation of expenses to departments in departmental accounting.\pts{14}

\medskip\rule{\linewidth}{0.2pt}




\noindent   \textbf{Question 5.}
Define Human Resource Accounting. Discuss its advantages and limitations and explain the `Present Value Method' for valuation of human resource.\pts{14}

\centerline{   \textbf{OR}}

Write an explanatory note on models of valuation of human resource.\pts{14}

\vfill
\rule{\linewidth}{0.4pt}

\begin{center}\small
\href{https://bhu-exam.vercel.app/}{Click here to visit our website}\\[4pt]
\href{https://me-aryan.vercel.app/}{Click here to visit our contributor}\\[4pt]
\href{https://quashan-me.vercel.app/}{Click here to visit our contributor}
\end{center}

\end{document}
